\section{Current Status and Reproduction}

\subsection{Where the project is now}

\begin{table}[H]
\centering
\caption{Publication-time status.}
\begin{tabularx}{\textwidth}{@{}L{0.23\textwidth}L{0.20\textwidth}Y@{}}
\toprule
Workstream & Status & Meaning \\
\midrule
Bar backtester & Educational precursor & Reusable framework and tests remain; historical performance claims are not treated as rigorous evidence. \\
Frozen L2 study & Complete under legacy model & Phase A/B/C artifacts, gates, panel identity, and audits are committed and verifiable. \\
Event-driven Python model & Mechanics accepted & Private scheduling, post-response-gated snapshot recovery, delayed cancels, own FIFO, conservation, gaps, and risk invariants pass deterministic tests. \\
V3 development economics & Pending & No execution-derived result from \CurrentModel\ is published yet. \\
Holdout & Strategy-sealed & Identities and regime descriptors were examined, but no candidate strategy run exists. \\
C++ performance port & Intentionally paused & It begins only after the Python V3 result is frozen and a separate modern C++ foundations phase is complete. \\
\bottomrule
\end{tabularx}
\end{table}

The immediate research milestone is the clean V3 development rerun under
\CurrentModel\ and \code{post_response_proxy_depth_boundary_v1}, including an explicit
before/after snapshot-timing comparison. It is not the C++ port and it is not a
holdout run.

\subsection{Fresh-clone verification}

The repository targets Python 3.11. At the 4 August 2026 checkpoint, a fresh
clone reports 376 deterministic tests passed and three optional raw-data smoke
checks skipped. All 379 pass when the selected captures are restored. From the
repository root:

\begin{lstlisting}[style=projectcode]
python3.11 -m venv .venv
source .venv/bin/activate
python -m pip install -r requirements.txt

env PYTHONPATH=. pytest -q
env PYTHONPATH=. python scripts/demo_event_driven_execution.py
env PYTHONPATH=. python scripts/verify_v2_artifacts.py
\end{lstlisting}

The synthetic simulator demo is deliberately tiny: a resting buy arrives, a
cancellation is requested, a recorded sell reaches the queue first, and the
delayed cancel is reported too late. It exercises execution lifecycle mechanics
without raw data; it does not exercise full replay timestamp batching, gap
recovery, empirical economics, or throughput.

The V2 verifier checks the manifest and panel identities, execution-model
marker, Phase A verdict, Phase B gate status, Phase C grid, same-millisecond
audit, panel-file hashes, and selected holdout identifiers across tracked
result paths and contents. It validates committed evidence; it does not
recompute replay results from raw messages.

\subsection{Building this report}

Report scalars and vector figures are regenerated only from an allowlist of
canonical frozen artifacts. The builder checks the execution-model marker and
input shape, including the distinction between 24 window rows and one pooled
row per queue/latency configuration.

\begin{lstlisting}[style=projectcode]
make report
make report-check
\end{lstlisting}

The first command rebuilds report inputs and the PDF. The second rebuilds and
then verifies that the generated asset manifest matches the committed copy.
Input/output SHA-256 values are in
\artifact{report/generated/asset_manifest.json}. The compiled document is
\artifact{report/l2_mm_research_report.pdf}.

\subsection{Full development rerun}

With the raw depth/trade layout described in \artifact{DATA.md}, the current
development runner defaults to the frozen development-window file and writes
only to the V3 namespace:

\begin{lstlisting}[style=projectcode]
env PYTHONPATH=. python scripts/run_l2_panel.py --phase all
\end{lstlisting}

Before relying on this command, use \code{--help} and inspect the committed
defaults. The runner verifies the development-panel file hash, refuses the
strategy-sealed holdout, and guards the frozen V2 directory. V3 artifact
verification remains impossible until this run produces its manifest.

\subsection{Suggested reviewer routes}

A reviewer can understand the project at three depths:

\begin{description}[leftmargin=6.2em,style=nextline]
  \item[Five minutes] Read the repository README, the executive summary, and
        the status table above.
  \item[Thirty minutes] Read Sections 3--7 and run the synthetic demo plus V2
        artifact verifier.
  \item[Code review] Start at \artifact{src/replay/engine.py} and
        \artifact{src/execution/simulator.py}; use
        \artifact{tests/test_execution_simulator.py} and
        \artifact{tests/test_engine.py} as executable specifications, then
        inspect \artifact{src/analysis/ofi_signal.py} for the research estimand.
\end{description}

\subsection{Roadmap to the performance boundary}

\begin{enumerate}
  \item Run the fixed baseline over all 24 development windows under
        \CurrentModel\ and the post-response-gated snapshot policy at both queue
        endpoints.
  \item Regenerate OFI conditional-on-fill evidence and queue/latency stress,
        then compare fill identity and economics with the legacy study.
  \item If a candidate clears a newly committed development rule, lock it and
        perform the one-shot holdout strategy evaluation; otherwise close the
        research arc without evaluating a candidate there.
  \item Freeze a deterministic Python reference, golden event traces, artifact
        schemas, and parity tolerances.
  \item Complete a separate modern C++ foundations phase before implementing
        any hot-path port. Later performance work must reproduce the Python
        contract before it is benchmarked.
\end{enumerate}

\enlargethispage{4\baselineskip}
\subsection{Conclusion}

The repository does not present a profitable market-making strategy. It
presents something more credible at its current stage: a tested microstructure
research system, a frozen negative experiment with auditable stopping logic,
and a documented correction from coarse private timing to an event-driven
execution reference. The remaining scientific question is whether the
development conclusions survive the current model. The project stops at
that boundary before any C++ performance implementation.
